\documentclass[../../../uem_tok_q.tex]{subfiles}

\begin{document}
    座標平面上を次の規則\ref{uem_tok_2024_f_sci_03_q:1}，%
    \ref{uem_tok_2024_f_sci_03_q:2}に従って1秒ごとに動く点Pを考える．

    \begin{enumerate}[label=(\roman*), align=right, itemindent=0pt,
        topsep=3pt, itemsep=3pt, leftmargin=3\zw, labelsep=.6\zw]
        \item 最初に，Pは点$(2,\;1)$にいる．\label{uem_tok_2024_f_sci_03_q:1}
        \item ある時刻でPが点$(a,\;b)$にいるとき，その1秒後にはPは
            
            \begin{itemize}[label={\textbullet}, leftmargin=2\zw, labelsep=.4\zw,
                topsep=2pt, itemsep=3pt]
                \item 確率 $\myfraca{1}{3}$ で$x$軸に関して$(a,\;b)$と対称な点
                \item 確率 $\myfraca{1}{3}$ で$y$軸に関して$(a,\;b)$と対称な点
                \item 確率 $\myfraca{1}{6}$ で直線 $y=x$ に関して$(a,\;b)$と対称な点
                \item 確率 $\myfraca{1}{6}$ で直線 $y=-x$ に関して$(a,\;b)$と対称な点
            \end{itemize}
            にいる．
            \label{uem_tok_2024_f_sci_03_q:2}
    \end{enumerate}
    以下の問いに答えよ．ただし，\ref{uem_tok_2024_f_sci_03_q:3}については，結論のみを書けばよい．
    \begin{enumerate}
        \item Pがとりうる点の座標をすべて求めよ．\label{uem_tok_2024_f_sci_03_q:3}
        \item $n$を正の整数とする．最初から$n$秒後にPが点$(2,\;1)$にいる確率と，%
            最初から$n$秒後にPが点$(-2,\;-1)$にいる確率は等しいことを示せ．
        \item $n$を正の整数とする．最初から$n$秒後にPが点$(2,\;1)$にいる確率を求めよ．
    \end{enumerate}
\end{document}
